\scanpage{334}
\mspage{73}
% archive sheet number: 333
% ms. 1
D'où une application sur les $\pi_0$, i.e.
% ms. 2
\[\bar f_{a,b}:\mathrm{Hom}_\Sigma(a,b)\longrightarrow\mathrm{Hom}_\Sigma(f(a),f(b)),\]
% ms. 3
donnée par
% ms. 4
\[\bar f_{a,b}(\mathrm{hot}(c))=\mathrm{hot}(f_{a,b}(c)).\]
% ms. 5–6
On a évidemment, pour deux chemins $\gamma:a\longrightarrow b$ et $\gamma':b\longrightarrow c$
% ms. 6
\[f_{a,\unsure{c|b}}(\gamma'\gamma)=f_{b,c}(\gamma')f_{a,b}(\gamma),\]
% ms. 7–12
ce qui prouve que $\bar f$ est compatible avec la composition. Il est trivial qu'il est compatible avec les identités (car $f_{a,a}$ transforme le chemin vide \inserted{en $a$} en le chemin vide en $f(a)$), donc on a un foncteur $\bar f$, et il rend bel et bien le carré commutatif.

% ms. 13–14
On peut maintenant prouver la propriété universelle de $\Pi_\Sigma(X)$ et de
% ms. 15
\[X\overset{\deleted{\unsure{p}}\unsure{k|p}}{\longleftrightarrow}\Pi_\Sigma(X),\]
% ms. 16–18
savoir que tout foncteur $f:X\longrightarrow X'$ tel que $u\in\Sigma\Longrightarrow f(u)\ \unsure{\mathrm{inv.}|\mathrm{iso.}}$, se factorise \uline{de façon unique} par $\Pi_\Sigma(X)$. L'unicité
% ms. 18–19 left
% slanted
\margin{C'est \unsure{inutile} \unsure{cf|cf.} ci-dessous}
% ms. 19–21
se voit comme tantôt\unsure{$*$|.} Pour l'existence, on note que $f$ \deleted{se factorise par $Y_{\unsure{is|1s}}$, la sous-catégorie pleine de $f$ ayant ses objets, et dont les} \inserted{applique $\Sigma$ dans l'ens. $\Sigma'$\unsure{$*$|:}}
